\section{Related Work}
\label{sec:related}

\subsection{Contextual access control}

RBAC represents permissions through roles and remains useful as a coarse
administrative mechanism~\cite{sandhu_rbac_1996}. It does not, by itself,
express whether an officer is currently assigned to a case, whether a request
crosses a jurisdiction, or whether the requested field is necessary for a
declared purpose. ABAC instead evaluates subject, object, action, and
environment attributes against policy~\cite{nist_abac_2014}; XACML is a
well-known policy language and architecture for this style of
evaluation~\cite{oasis_xacml_2013}. Work on Fabric-based ABAC and accountable
privacy-preserving access control demonstrates that fine-grained policy and
ledger-backed accountability can be combined~\cite{zhao_fabric_abac_2022,hu_accountable_private_ac_2024}.

SEBA-XAI uses this lineage but makes a narrower systems choice. Subject
attributes are obtained from the signing certificate and current ledger state,
not accepted as request fields. Object attributes are read from committed
metadata. The requester supplies only an allow-listed action and purpose. The
contribution is therefore not ABAC itself, but an implementation in which the
authorisation result and a structured justification arise from the same
endorsed transaction.

\subsection{Permissioned audit and evidence provenance}

Permissioned ledgers are suited to organisations that need known membership,
replicated history, and policy-governed writes. Hyperledger Fabric provides the
relevant building blocks: membership identities, endorsing peers, ordering, and
application chaincode~\cite{androulaki_fabric_2018}. NIST's analysis of
blockchain access control likewise identifies auditability and shared
governance as potential benefits while retaining privacy and scalability as
design constraints~\cite{nist_blockchain_access_2022}.

Prior work applies blockchain to digital-evidence lifecycle and custody
management~\cite{jeong_fabric_evidence_2020,kim_two_level_2021,li_lechain_2021}
and to auditable distributed business-process access
control~\cite{akhtar_auditable_access_2020}. These designs motivate recording
commitments rather than sensitive payloads. They do not imply that a ledger
entry is correct merely because it is immutable. A ledger can establish
agreement about a committed transaction; it cannot independently establish the
truth of compromised identity attributes or other premises used during policy
evaluation. This distinction drives the threat analysis in
Section~\ref{sec:method}.

\subsection{Explainability for high-stakes review}

Explanation methods such as LIME and SHAP are often used to interpret learned
models~\cite{ribeiro_lime_2016,lundberg_shap_2017}, and counterfactuals offer a
way to describe changes that would alter an automated
outcome~\cite{wachter_counterfactual_2018}. For a deterministic access policy,
however, a rule trace is more direct than a post-hoc attribution. This choice
also follows the caution that high-stakes decisions should prefer inherently
interpretable mechanisms where they are adequate~\cite{rudin_2019}. Recent
law-enforcement XAI work stresses stakeholder-specific explanations and the
risk of automation bias~\cite{zocholl_xai_law_enforcement_2025}.

Blockchain-supported XAI audit proposals have considered anchoring explanation
artifacts~\cite{nassar_blockchain_xai_2020,malhotra_xai_audit_2021}. SEBA-XAI
uses no model to make an access decision. Instead, it records the deterministic
reason code, decisive attributes, policy version, and counterfactual as a
machine-checkable explanation artifact that is hash-linked to the audit event.
Natural-language rendering, when enabled in the prototype, is a post-decision
presentation step and cannot affect authorisation.

